\chapter{Bounded curvature and corrected mass ratios}\label{sec:varifold-bounded-curvature}
Here $V$ is an $m$-varifold in $U\subset\mathbb R^N$, $1\leq m<N$, with weight $\mu$. Assume throughout that
\[
 \delta V(X)=-\int H\cdot X\,d\mu,\qquad |H|\leq\Lambda\quad\mu\text{-almost everywhere},
\]
for compactly supported $C^1$ fields, with $\Lambda\geq0$. In particular the first variation has no singular boundary part. These are sufficient hypotheses for \cite{Sim18}, Chapter 8, (3.4)--(3.5), p.~241. No rectifiability is required unless stated.

\begin{proposition}[Annular identity with mean-curvature error]\label{prop:curvature-annular-identity}
Let $\overline{B_R(a)}\subset U$ and $0<\sigma<\rho<R$. With $f_a,T_a$ as in Definition \ref{def:radial-varifold-tilt}, define
\[
 \gamma_{\sigma,\rho}(r)=\begin{cases}
 \max\{r,\sigma\}^{-m}-\rho^{-m},&r<\rho,\\0,&r\geq\rho.
 \end{cases}
\]
Then
\[
 f_a(\rho)-f_a(\sigma)
 =\int_{G_m(B_\rho(a)\setminus B_\sigma(a))}T_a\,dV
 +\frac1m\int\gamma_{\sigma,\rho}(|x-a|)H(x)\cdot(x-a)\,d\mu.
\]
In particular
\[
 f_a(\rho)-f_a(\sigma)\geq-\Lambda\int_\sigma^\rho f_a(t)\,dt.
\]
\end{proposition}
\begin{proof}
Use the cutoff and divergence computation in the proof of Proposition \ref{prop:stationary-annular-identity}. First variation is now $-\int\gamma H\cdot(x-a)\,d\mu$, giving the stated sign after rearrangement. Approximate arbitrary radii from below through continuity radii as before; the bounded curvature and local finiteness justify dominated convergence for the extra term.

The layer-cake formula
\[
 \gamma_{\sigma,\rho}(r)=\int_\sigma^\rho m t^{-m-1}\boldsymbol1_{\{r<t\}}\,dt
\]
shows, by Tonelli, that the magnitude of a negative curvature contribution is at most
\[
 \Lambda\int_\sigma^\rho t^{-m-1}
       \int_{B_t(a)}|x-a|\,d\mu\,dt
 \leq\Lambda\int_\sigma^\rho f_a(t)\,dt.
\]
The radial-tilt term is nonnegative. No differentiability of the ball-mass function was used.
\end{proof}

\begin{lemma}[An integrating factor for a measure differential inequality]\label{lem:mass-ratio-integrating-factor}
Let $f$ be a locally bounded, left-continuous function of locally bounded variation on an interval $I$. If for every $s<t$ in $I$,
\[
 f(t)-f(s)\geq-\Lambda\int_s^t f(r)\,dr,
\]
then $e^{\Lambda t}f(t)$ is nondecreasing.
\end{lemma}
\begin{proof}
Fix $t_0\in I$ and put $g(t)=f(t)+\Lambda\int_{t_0}^t f(r)\,dr$. The hypothesis says $g$ is nondecreasing, so its distributional derivative is a nonnegative Radon measure. The product rule for a smooth factor and a BV function gives
\[
 D(e^{\Lambda t}f)=e^{\Lambda t}(Df+\Lambda f\,dt)
                  =e^{\Lambda t}Dg\geq0.
\]
For completeness this product rule follows by writing the distributional derivative against a compactly supported smooth test and expanding the derivative of its product with $e^{\Lambda t}$. A function with nonnegative distributional derivative has a nondecreasing representative, as is seen by mollification. Left continuity selects that representative at every interior point, including jumps, and proves the claimed pointwise conclusion.
\end{proof}

\begin{theorem}[Exponential monotonicity for bounded mean curvature]\label{thm:bounded-curvature-monotonicity}
Let $V$ be an $m$-varifold in $U\subset\mathbb R^N$, $1\leq m<N$, with $\mu=\|V\|$. Suppose $\delta V(X)=-\int H\cdot X\,d\mu$ for all $X\in C_c^1(U;\mathbb R^N)$ and $|H|\leq\Lambda$ almost everywhere, where $\Lambda\geq0$. If $\overline{B_R(a)}\subset U$, then
\[
 r\longmapsto e^{\Lambda r}\frac{\mu(B_r(a))}{\omega_m r^m}
\]
is nondecreasing for $0<r<R$. The density $\Theta^m(\mu,a)$ exists, is finite and nonnegative, and $\mu(\{a\})=0$. This is the bounded-mean-curvature case of \cite{Sim18}, Chapter 8, (3.5), p.~241.
\end{theorem}
\begin{proof}
The function $r\mapsto\mu(B_r(a))$ is nondecreasing and left-continuous. Therefore $f_a$ is locally bounded and of locally bounded variation away from zero. Proposition \ref{prop:curvature-annular-identity} and Lemma \ref{lem:mass-ratio-integrating-factor} give exponential monotonicity. Its nonnegative ratios have a finite limit at zero, bounded above by any fixed-radius value. Since $e^{\Lambda r}\to1$, that limit is the usual density after division by $\omega_m$. The resulting small-ball bound $\mu(B_r(a))\leq Cr^m$ excludes atoms.
\end{proof}

\begin{corollary}[Density and mass lower bounds with bounded curvature]\label{cor:bounded-curvature-support-density}
If $V$ is also integral, then for every $a\in\operatorname{spt}\mu$,
\[
 \Theta^m(\mu,a)\geq1,\qquad
 \mu(B_r(a))\geq e^{-\Lambda r}\omega_m r^m
\]
whenever $\overline{B_r(a)}\subset U$.
\end{corollary}
\begin{proof}
First density is upper semicontinuous in the center. Indeed, for $a_j\to a$ and $0<\varepsilon<s$, the inclusion $B_{s-\varepsilon}(a_j)\subset B_s(a)$ and exponential monotonicity give
\[
 \Theta^m(\mu,a_j)\leq
 e^{\Lambda(s-\varepsilon)}\frac{\mu(B_s(a))}{\omega_m(s-\varepsilon)^m}
\]
for large $j$. Let $\varepsilon\downarrow0$, then $s\downarrow0$. Integrality supplies density at least one almost everywhere. Choose such points approaching any support point; upper semicontinuity gives the same lower bound there. Finally compare the corrected ratio at $r$ with its limit at zero. This proof parallels Corollary \ref{cor:integral-support-density} and specifies the curvature correction.
\end{proof}

\begin{example}[Why a disk is not a zero-curvature test case]\label{ex:disk-singular-first-variation}
A flat unit $m$-disk in an $m$-plane of $\mathbb R^N$, viewed as a varifold on all of $\mathbb R^N$, has a boundary contribution to first variation. It does not satisfy the hypotheses of this section with $H=0$.
\end{example}
\begin{proof}
The divergence theorem gives $\delta V(X)=\int_{\partial D}X\cdot\nu\,d\mathcal H^{m-1}$, where $\nu$ is the outward conormal; for $m=1$ this is the signed endpoint sum. This measure is singular with respect to the disk's $m$-dimensional weight. The centered ratio $\|V\|(B_r)/r^m$ equals $\omega_m$ for $r\leq1$ and $\omega_m r^{-m}$ for $r>1$, so it decreases beyond the boundary. Vanishing classical curvature in the interior cannot remove that boundary term.
\end{proof}
